\section{Finite-lattice and heat-kernel conventions}\label{sec:setup}

\subsection{The open hypercubic block}

Fix $d\geq2$.  Let $B_d$ be the oriented cubical complex whose vertices are
\[
  B_d^{(0)}=\{0,1,2\}^d.
\]
Every positively oriented nearest-neighbor edge is a link and every positively
oriented unit coordinate square is a plaquette.  A gauge configuration is
$U=(U_a)_{a\in B_d^{(1)}}\in G^{B_d^{(1)}}$, with
$U_{\bar a}=U_a^{-1}$ on the reversed orientation.  For a plaquette $p$ with
oriented boundary $(a_1,a_2,a_3,a_4)$, its holonomy is
\[
  U_p=U_{a_1}U_{a_2}U_{a_3}U_{a_4}.
\]
Changing the initial vertex conjugates $U_p$, so its distance from the
identity is unambiguous when the metric is bi-invariant.

Set
\[
  x=(0,1,\ldots,1),\qquad e=(x,1),
\]
where coordinate directions are numbered $1,\ldots,d$.  The link $e$ runs
from $x$ to $x+\mathbf e_1$.  For every $\nu=2,\ldots,d$, it lies in the two
unit squares extending from transverse coordinate $1$ toward $0$ and toward
$2$.  Consequently its incident-plaquette degree is
\begin{equation}\label{eq:m-definition}
  m=2(d-1).
\end{equation}
We call such a link \emph{full-degree}; no claim that it is disjoint from the
geometric boundary in its own direction is intended.

\subsection{Compact group and heat kernel}

Let $G$ be a compact connected Lie group equipped with a bi-invariant
Riemannian metric.  Denote its normalized Haar probability measure by
$\Haar$, its identity by $\one$, and its Riemannian distance by $\dist$.
The heat kernel of $\tfrac12\Delta$ based at $\one$ is written $p_t(U)$.
Thus $p_t$ is strictly positive, smooth, central, invariant under inversion,
and normalized by
\[
  \int_G p_t(U)\,\Haar=1.
\]
The heat-kernel plaquette measure on $B_d$ is
\begin{equation}\label{eq:finite-measure}
 \mu_t(\dd U)=\mathcal Z_t^{-1}
 \prod_{p\in B_d^{(2)}}p_t(U_p)
 \prod_{a\in B_d^{(1)}}\dd U_a.
\end{equation}
This is a finite-dimensional probability measure with a strictly positive
smooth density.  The heat-kernel action in lattice gauge theory goes back to
Menotti and Onofri \cite{MenottiOnofri1981}.  We use it because its
short-time geometry is exact at logarithmic scale.

For each fixed link $e$, the positivity and continuity of
\eqref{eq:finite-measure} give a canonical continuous version of the
one-link conditional kernel for every assignment
$\omega\in G^{B_d^{(1)}\setminus\{e\}}$.  No almost-everywhere ambiguity is
therefore involved when we evaluate the kernel at a center-valued boundary.

\subsection{Staples and defect events}

After cyclically reordering a plaquette product and, when necessary, using
$p_t(V)=p_t(V^{-1})$, each incident factor can be written
\begin{equation}\label{eq:staple-form}
  p_t(U_e S_p(\omega))
\end{equation}
for a continuous staple map
$S_p:G^{B_d^{(1)}\setminus\{e\}}\to G$.  The precise orientation convention
for $S_p$ is fixed by \eqref{eq:staple-form}; a different convention only
inverts the prescribed central value.

For an incident target plaquette $q$ and $0<\rho<\operatorname{diam}G$, set
\begin{equation}\label{eq:good-event}
 \mathcal G_{q,\rho}(\omega)
 =\{U\in G:\dist(\one,U S_q(\omega))<\rho\}.
\end{equation}
Its complement is the local bad-curvature event.  The boundary-uniform
rare-defect estimate addressed here is
\begin{equation}\label{eq:uniform-quantifier}
 \sup_{\omega\in G^{B_d^{(1)}\setminus\{e\}}}
 \gamma_{e,t}^{\omega}(\mathcal G_{q,\rho}(\omega)^c)
 \leq \varepsilon(t),
 \qquad \varepsilon(t)\xrightarrow[t\downarrow0]{}0.
\end{equation}
The event is gauge invariant, although the link-coordinate conditioning is
not.  Theorem~\ref{thm:central-obstruction} proves that the left side of
\eqref{eq:uniform-quantifier} instead tends to one.

\subsection{Dobrushin convention}

For distinct links $a,b$, define the finite-volume interdependence
coefficient
\begin{equation}\label{eq:dobrushin-coefficient}
 c_{ab}(t)=
 \sup_{\substack{\omega,\omega'\\
                  \omega_c=\omega'_c\ (c\neq b)}}
 \left\|\gamma_{a,t}^{\omega}-\gamma_{a,t}^{\omega'}\right\|_{\TV},
 \qquad
 \|\nu-\nu'\|_{\TV}=\sup_A|\nu(A)-\nu'(A)|.
\end{equation}
The classical row-sum condition is
\begin{equation}\label{eq:dobrushin-row-condition}
  \max_a\sum_{b\neq a}c_{ab}(t)<1.
\end{equation}
It is sufficient for standard uniqueness and comparison conclusions in the
appropriate infinite-volume setting \cite{Dobrushin1968}.  Showing that
\eqref{eq:dobrushin-row-condition} fails does not prove nonuniqueness.
